\subsection{函数空间在施瓦茨函数空间中的包含}

命题 12. — 空间 \(\mathcal{D}(\mathbf{R}^{n})\) 含于 \(\mathcal{S}(\mathbf{R}^{n})\) ，且 \(\mathcal{D}(\mathbf{R}^{n})\) 到 \(\mathcal{S}(\mathbf{R}^{n})\) 中的包含是连续的且像稠密。

设 \(B \subset \mathcal{D}(\mathbf{R}^{n})\) 是有界子集，K 是 \(\mathbf{R}^{n}\) 的紧子集且 \(\mathrm{B} \subset \mathcal{C}_{\mathrm{K}}^{\infty}(\mathbf{R}^{n})\) 。设 \(\alpha \in \mathbf{N}^{n}\) 与 \(k \in \mathbf{N}\) 。对每个函数 \(\varphi \in B\) ，可得

\[
\widetilde {q} _ {\alpha , k} (\varphi) \leqslant \Bigl (\sup _ {x \in \mathrm{K}} \| x \| ^ {k} \Bigr) p _ {\alpha , \mathrm{K}} (\varphi),
\]

因而 B 在 \(\mathcal{S}(\mathbf{R}^{n})\) 中有界。包含的连续性则由空间 \(\mathcal{S}(\mathbf{R}^{n})\) 与 \(\mathcal{D}(\mathbf{R}^{n})\) 都是有界型的这一事实得出。

我们来证明 \(\mathcal{D}(\mathbf{R}^n)\) 在 \(\mathcal{S}(\mathbf{R}^n)\) 中稠密。设 B 是 \(\mathbf{R}^n\) 的单位球，\(\eta \in \mathcal{D}(\mathbf{R}^n)\) 是支集含于 2B 且满足 \(0 \leqslant \eta \leqslant 1\) 及 \(\eta(x) = 1\) （对所有 \(x \in \mathrm{B}\) ）的测试函数（IV, p. 196，引理 1, \(a\) ）。

设 \(\varphi \in \mathcal{S}(\mathbf{R}^n)\) 。对每个整数 \(m \geqslant 1\) 及每个 \(x \in \mathbf{R}^n\) ，令 \(\eta_m(x) = \eta(x/m)\) 。最后定义 \(\varphi_m = \eta_m\varphi\) ；我们有 \(\varphi_m \in \mathcal{D}(\mathbf{R}^n)\) 。

由于 \(\partial^{\alpha}\eta_{m} = m^{-|\alpha|}(\partial^{\alpha}\eta)(x / m)\) （对所有 \(\alpha \in \mathbf{N}^{n}\) 与 \(x\in \mathbf{R}^n\) ），由 IV, p. 210 的公式 (2) 推出序列 \((\varphi_{m})\) 在 \(\mathcal{S}(\mathbf{R}^n)\) 中有界。

设 C 是 \(R^{n}\) 的紧子集。序列 \((\varphi_{m})_{m\geqslant1}\) 收敛于 \(\varphi\) （在 \(\mathcal{C}_{\mathrm{K}}^{\infty}(\mathbf{R}^{n})\) 中），因为对所有充分大的 m，\(\varphi_{m}\) 在 C 上与 \(\varphi\) 一致。于是序列 \((\varphi_{m})\) 收敛于 \(\varphi\) （在 \(\mathcal{C}^{\infty}(\mathbf{R}^{n})\) 中），而 IV, p. 211 的命题 10 的推论使我们得出序列 \((\varphi_{m})\) 收敛于 \(\varphi\) （在 \(\mathcal{S}(\mathbf{R}^{n})\) 中）。

引理 8. --- 设 B 是 \(\mathcal{D}(\mathbf{R}^{n})\) 的有界子集。\(\mathcal{S}(\mathbf{R}^{n})\) 的拓扑在 B 上诱导的拓扑与 \(\mathcal{D}(\mathbf{R}^{n})\) 诱导的拓扑一致。

由于 \(\mathcal{D}(\mathbf{R}^{n})\) 到 \(\mathcal{S}(\mathbf{R}^{n})\) 中的包含是连续的，\(\mathcal{D}(\mathbf{R}^{n})\) 在 B 上诱导的拓扑比 \(\mathcal{S}(\mathbf{R}^{n})\) 诱导的拓扑更细。此外，存在 \(\mathbf{R}^{n}\) 的紧子集 K 使得 \(\mathrm{B} \subset \mathcal{C}_{\mathrm{K}}^{\infty}(\mathbf{R}^{n})\) 。于是对每个 \(\alpha \in \mathbf{N}^{n}\) ，有 \(p_{\alpha,K}(\varphi) \leqslant \widetilde{q}_{\alpha,0}(\varphi)\) ，这蕴含 \(\mathcal{S}(\mathbf{R}^{n})\) 诱导的拓扑比 \(\mathcal{D}(\mathbf{R}^{n})\) 诱导的拓扑更细。

命题 13. --- 设 \(p \in [1, +\infty]\) 。空间 \(\mathcal{S}(\mathbf{R}^{n})\) 含于 \(\mathcal{L}^{p}(\mathbf{R}^{n})\) ，且 \(\mathcal{S}(\mathbf{R}^{n})\) 到 \(\mathcal{L}^{p}(\mathbf{R}^{n})\) 中的典范单射是连续的。则 \(\mathcal{S}(\mathbf{R}^{n})\) 在 \(\mathrm{L}^{p}(\mathbf{R}^{n})\) 中的像（若 \(p \neq +\infty\) ）稠密。

对 \(p = +\infty\) ，第一个断言是直接的。现设 \(p \in [1, +\infty[\) 。设 m 是满足 \(n + 1 < mp\) 的整数。对每个 \(\varphi \in \mathcal{S}(\mathbf{R}^{n})\) 与 \(x \in \mathbf{R}^{n}\) ，我们有

\[
\| x \| ^ {n + 1} | \varphi (x) | ^ {p} \leqslant \widetilde {q} _ {0, m} (\varphi) ^ {p}
\]

因而依据 IV, p. 199 的命题 3，\(\varphi \in \mathcal{L}^p (\mathbf{R}^n)\) 。此外，我们得到

\[
\mathrm{N} _ {p} (\varphi) \leqslant a _ {n} ^ {1 / p} \widetilde {q} _ {0, 0} (\varphi) + b _ {n} \widetilde {q} _ {0, m} (\varphi),
\]

其中

\[
a _ {n} = \int_ {\| x \| \leqslant 1} d \mu (x), \quad b _ {n} = \int_ {\| x \| \geqslant 1} \frac {1}{\| x \| ^ {n + 1}} d \mu (x),
\]

因此 \(\mathcal{S}(\mathbf{R}^n)\) 到 \(\mathcal{L}^p (\mathbf{R}^n)\) 中的单射是连续的。

由于空间 \(\mathcal{D}(\mathbf{R}^n)\) 含于 \(\mathcal{S}(\mathbf{R}^n)\) ，IV, p. 202 的命题 4 蕴含 \(\mathcal{S}(\mathbf{R}^n)\) 在 \(\mathrm{L}^p (\mathbf{R}^n)\) 中稠密（若 \(p\neq +\infty\) ）。

